
\begin{claim}{AC.34}{Proved; limitation of a trigonal branch-point argument}
Let \(E/L\) be an ordinary elliptic curve in characteristic two,
and let \(f:E\to\mathbf P^1\) be a separable degree-three map
whose ramification consists of three distinct simple wild points
\(R_1,R_2,R_3\). No difference \(R_i-R_j\), \(i\ne j\), is the
nonzero geometric two-torsion class. It is nevertheless possible
for all three differences to be two-primary.

An explicit example over \(\mathbf F_4\), with
\(\alpha^2+\alpha+1=0\) and \(\beta=\alpha^2\), is
\[
 E:\ y^2+xy=x^3+1,\qquad f=(y+\alpha)/x.
\]
The ramification points are
\[
 T=(0,1),\quad R=(\beta,0),\quad R'=(\beta,\beta),
\]
and the pair-difference orders are \(8,8,4\).
The unramified residual points in their degree-three fibres are
\[
 S_\infty=\infty,\qquad S_\beta=(1,1),\qquad S_\alpha=(1,0),
\]
whose pair-difference orders are \(4,4,2\).
The quadratic resolvent is the genus-two curve
\[
 C:\quad y^2+y=z+\frac{\beta}{z^2+z}+\alpha.
\]
Its constant pointed cubic cover has
\(L(\chi,T)=1+3T+4T^2\) for both nontrivial characters and
relative class number \(64\).

For a pointed dihedral cubic cover \(D\to C\), twisting by the
reflection fixing the chosen lift leaves its elliptic quotient
unchanged. Thus this twist cannot reduce the two-primary order
of the elliptic group in the proposed argument.
These statements exhibit a limitation of that route, not a
collision on \(C\).
\end{claim}
\begin{proof}
Let \(\eta\ne0\) be the unique geometric two-torsion point of
\(E\). The locally exact differential line on an ordinary
elliptic curve is the nontrivial two-torsion line bundle
corresponding to \(\eta\).
If \(D_\infty=f^*(\infty)\), the ramification formula gives
\[
 \operatorname{div}(df)
        =2(R_1+R_2+R_3-D_\infty),\qquad
 [R_1+R_2+R_3-D_\infty]=\eta.
\]
For completeness the class of half this exact differential
cannot be zero: otherwise \(df=g^2\omega\) for a nonzero
regular differential \(\omega\); applying Cartier gives
\(0=g\,\mathcal C(\omega)\), contrary to ordinariness.
If \(R_1-R_2=\eta\), cancellation would give
\(D_\infty\sim2R_2+R_3\).
The actual fibre at \(f(R_2)\) is \(2R_2+S_2\).
Cancelling \(2R_2\) and using the injective Abel map of \(E\)
gives \(S_2=R_3\). That fibre would contain both ramification
points with multiplicity at least two, contradicting degree
three. This proves the order-two assertion.

For the example, direct enumeration gives
\(E(\mathbf F_4)\simeq\mathbf Z/8\).
The pole divisor of \(f\) is \(2T+\infty\), so its degree is
three, and with \(\omega=dx/x\) one has
\[
                        df=\frac{x^2+\alpha}{x}\,\omega.
\]
The divisor of this differential and the pole divisor give
different exponent two at exactly \(T,R,R'\).
Their branch values are respectively \(\infty,\beta,\alpha\).
The stated residual points follow by evaluating \(f\);
exact elliptic addition gives the six difference orders in
the claim. In particular the distinction between order two
and a higher two-power is essential.

Writing \(f=t\), so \(y=tx+\alpha\), gives the cubic
\[
                x^3+(t^2+t)x^2+\alpha x+\alpha=0.
\]
For a separable cubic \(X^3+aX^2+bX+c\) in characteristic two,
its quadratic resolvent has Artin--Schreier class
\((b^3+a^3c+c^2)/(ab+c)^2\).
Here this simplifies to
\[
              \beta(t^2+t)+\frac{\beta}{t^2+t+1}.
\]
The change \(w\mapsto w+\alpha t\), followed by \(t=z+\alpha\),
gives the displayed \(C\).
The Galois closure has group \(S_3\), since the cubic is
irreducible over \(L(t)\) by its degree-three function-field
map and the nontrivial quadratic resolvent is connected.
All three branch inertia groups have order two, so the
degree-three map from this closure to \(C\) is unramified.
The rational residual points give rational lifts of the three
branch points of \(C\), making this constant cubic cover
pointed. Alternatively its unique pointed character is
identified directly in \(J_C(\mathbf F_4)\), which has order
twelve. The three rational points of \(C\) are its branch
points and have character occupancies \((3,0,0)\).
AC.4 therefore gives \(1+3T+4T^2\) and relative value
\(8^2=64\). All the finite-field, group, resolvent and
character calculations are executed in
\path{trigonal_two_primary_obstruction.py}.

Finally, if \(D\to C\) has deck group \(\langle\delta\rangle\)
of order three and \(\tau\) is its pointed reflection,
then a \(\tau\)-fixed lift over a branch point maps to the
unramified residual point \(S_i\) of
\(E=D/\langle\tau\rangle\to\mathbf P^1\), not to \(R_i\).
Indeed the total inertia has order two and is already
accounted for in \(D\to E\). The fibre relation is
\[
                     2R_i+S_i\sim D_\infty,\qquad
                     S_i-S_j=-2(R_i-R_j).
\]
Twisting \(D\) by \(\tau\) does not twist \(E\), because
\(\tau\) acts trivially on its quotient. Every other reflection
\(\delta^j\tau\) is conjugate to \(\tau\) by a rational deck
element, so it gives the same elliptic quotient up to
\(L\)-isomorphism when the deck group is constant.
A nontrivial pure deck twist destroys the rational lift over
the prescribed branch point. Thus neither operation supplies
the missing small-two-primary-group hypothesis.
\end{proof}

